\clearpage
\section{Pairwise Berry connection
comparison}\label{f-pairwise-berry-connection-comparison}

Technical derivation for review. Equation and subsection numbers in this
appendix are local to the source derivation. Historical local labels are
retained, including numbering gaps or nonsequential labels. The
accompanying source notes retain the original expressions for
comparison.

\subsection{4.1 Complete charge
comparison}\label{f-complete-charge-comparison}

A naive kinetic comparison

% DISPLAY f.1
% SOURCE |(partial+iA)f|² >= (1−epsilon)|partial f|²−epsilon^(-1)|A|²|f|²
\begin{gather*}
\begin{aligned}
& \vert (\partial +iA)f\vert ^{2} \ge  (1-\epsilon )\vert \partial  f\vert ^{2}-\epsilon ^{-1}\vert A\vert ^{2}\vert f\vert ^{2}
\end{aligned}\notag
\end{gather*}

loses a fixed part of internal kinetic energy. Replacing it by \(h_{I}\)
incurs the unwanted \(O(\epsilon  \kappa  r)\) internal Yukawa mass.
This is the wrong comparison.

Instead, normalize the sixteen internal charges so \(h_{I}\) is their
averaged squared norm on the core. Replace each momentum by its scalar
Berry-covariant momentum. Direct Clifford expansion gives

% DISPLAY f.2
% SOURCE h_I(A)= average_alpha ||Q_(I,alpha)+C_alpha(A)||²
% SOURCE          − a finite Clifford contraction of F_A.
\begin{gather*}
\begin{aligned}
& h_{I}(A)= \operatorname{average} _{\alpha} \Vert Q_{I,\alpha}+C_{\alpha}(A)\Vert ^{2} \\
& - \text{ a finite Clifford contraction of } F_{A}.
\end{aligned}\notag
\end{gather*}

The difference \(C_{\alpha}(A)\) is a bounded multiplication matrix,
with norm \(\le C_{N}\vert A\vert\). Consequently

% DISPLAY f.3
% SOURCE h_I(A) >= (1−epsilon) h_I
% SOURCE            −C_N[epsilon^(-1)|A|²+|F_A|].
\begin{gather*}
\begin{aligned}
& h_{I}(A) \ge  (1-\epsilon ) h_{I} \\
& -C_{N}[\epsilon ^{-1}\vert A\vert ^{2}+\vert F_{A}\vert ].
\end{aligned}\notag
\end{gather*}

This retains the complete interacting \(h_{I}\), rather than its kinetic
part, and uses no gap. Spinor averaging eliminates the ordinary
gauge-term obstruction to positivity even on the unrestricted internal
module; the low coefficient additionally has the physical residual
equivariance required by the inductive Hardy input.

On a two-block small-ratio tube the gapped finite matrix is analytic.
All determinant curvatures vanish at internal \(A=0\): the two-projector
derivative calculation contains a Spin(9) trace of at most three
Clifford matrices. Homogeneity and smoothness therefore give

% DISPLAY f.4
% SOURCE |F_A| <= C_N kappa/r².
\begin{gather*}
\begin{aligned}
& \vert F_{A}\vert  \le  C_{N} \kappa /r^{2}.
\end{aligned}\notag
\end{gather*}

Use equivariant radial gauge in the internal coordinates, based on the
flat center vacuum at \(A=0\). The radial-gauge curvature formula gives
\(\vert A_{A}\vert \le C_{N} \kappa ^{2}/r\) for internal components.
Mixed relative components must also be kept: in this gauge they are the
flat \(A=0\) base connection plus the integral of \(F_{A,b}\) contracted
with A. The same vanishing-at-\(A=0\) argument and homogeneity give
their bound \(C_{N} \kappa ^{2}/r\). Thus, after this separate
mixed-component check,

% DISPLAY f.5
% SOURCE |A_Berry| <= C_N kappa²/r.
\begin{gather*}
\begin{aligned}
& \vert A_{\mathrm{Berry}}\vert  \le  C_{N} \kappa ^{2}/r.
\end{aligned}\notag
\end{gather*}

It follows that the preceding comparison costs only

% DISPLAY f.6
% SOURCE C_N[kappa+epsilon^(-1) kappa^4] r^(-2),
\begin{gather*}
\begin{aligned}
& C_{N}[\kappa +\epsilon ^{-1} \kappa ^{4}] r^{-2},
\end{aligned}\notag
\end{gather*}

which can be made arbitrarily small after fixing \(\epsilon\). A
gauge-equivariant radial trivialization preserves ordinary residual
SU(n) physicality of the coefficient. Weyl/center flat identifications
remain part of the bundle.

This establishes a viable repair of the flatness step provided the local
Hamiltonian comparison really identifies the covariant internal form and
retains its curvature terms. The following rootwise construction avoids
an artificial dependence on the global minimum-gap ratio.

\subsection{4.2 Pairwise relative multiblock Berry
bound}\label{f-pairwise-relative-multiblock-berry-bound}

At quadratic order the fast fermionic space is a direct sum over
bifundamental block pairs. Its one-particle matrices are

% DISPLAY f.7
% SOURCE D_ab=Gamma dot b_ab tensor I
% SOURCE        +sum_i Gamma_i tensor (A_i^(a) tensor I−I tensor (A_i^(b))^T).
\begin{gather*}
\begin{aligned}
& D_{ab}=\Gamma  \cdot  b_{ab} \otimes  I \\
& +\sum _{i} \Gamma _{i} \otimes  (A_{i}^{(a)} \otimes  I-I \otimes  (A_{i}^{(b)})^{T}).
\end{aligned}\notag
\end{gather*}

If each endpoint internal norm is at most \(\kappa\) times
\(r_{ab}=\vert b_{ab}\vert\), the gap is a positive fixed multiple of
\(r_{ab}\) after reducing \(\kappa\) by the fixed nine-dimensional norm
factor. The adapted determinant line is the tensor product of the
occupied pair determinants. On each pair, choose the preceding
internal-radial gauge from the flat \(A=0\) center line. The Clifford
trace calculation at \(A=0\) is independent of bifundamental dimension,
so every pure/mixed curvature component vanishes there.

Smoothness on the compact dimensionless pair tube and homogeneity give,
in canonical block/center coordinates,

% DISPLAY f.8
% SOURCE |F_ab|<=C_N kappa r_ab^(-2),
% SOURCE |a_ab|<=C_N kappa² r_ab^(-1).
\begin{gather*}
\begin{aligned}
& \vert F_{ab}\vert \le C_{N} \kappa  r_{ab}^{-2}, \\
& \vert a_{ab}\vert \le C_{N} \kappa ^{2} r_{ab}^{-1}.
\end{aligned}\notag
\end{gather*}

The relative-component bound follows from its own mixed-curvature
integral, not from simply declaring it gauged away. Product connections
add, and Cauchy--Schwarz over the finitely many pairs gives

% DISPLAY f.9
% SOURCE |F_total|<=C_N kappa sum_(a<b) r_ab^(-2),
% SOURCE |a_total|²<=C_N kappa^4 sum_(a<b) r_ab^(-2).
\begin{gather*}
\begin{aligned}
& \vert F_{\mathrm{total}}\vert \le C_{N} \kappa  \sum _{a<b} r_{ab}^{-2}, \\
& \vert a_{\mathrm{total}}\vert ^{2}\le C_{N} \kappa ^{4} \sum _{a<b} r_{ab}^{-2}.
\end{aligned}\notag
\end{gather*}

The full slow averaged charge stack includes all internal interacting
charges and free center charges. Its comparison therefore loses only
\(\epsilon\) times the entire slow supersymmetric form and the
inverse-square multiplier

% DISPLAY f.10
% SOURCE C_N(kappa+epsilon^(-1)kappa^4) sum_(a<b) r_ab^(-2).
\begin{gather*}
\begin{aligned}
& C_{N}(\kappa +\epsilon ^{-1}\kappa ^{4}) \sum _{a<b} r_{ab}^{-2}.
\end{aligned}\notag
\end{gather*}

This multiplier has a direct shape-independent center-kinetic payment.
In the mass metric \(\sum  n_{a}\vert dz_{a}\vert ^{2}\), the linear map
\(z\to z_{a}-z_{b}\) has squared norm \(l_{ab}^{2}=1/n_{a}+1/n_{b}\).
Let \(T_{ab}\) be kinetic energy in its nine-dimensional orthogonal
relative subspace. Then

% DISPLAY f.11
% SOURCE T_ab >= (49/4) l_ab² integral r_ab^(-2)|f|²,
% SOURCE T_ab <= T_centers.
\begin{gather*}
\begin{aligned}
& T_{ab} \ge  (49/4) l_{ab}^{2} \int  r_{ab}^{-2}\vert f\vert ^{2}, \\
& T_{ab} \le  T_{\mathrm{centers}}.
\end{aligned}\notag
\end{gather*}

Thus sum
\(\int  r_{ab}^{-2}\vert f\vert ^{2} \le (4/49) \sum  l_{ab}^{-2} T_{\mathrm{centers}}\).
The finite rank-dependent coefficient contains no inverse shape gap
\(\delta _{N}\). Choose \(\epsilon\), then \(\kappa\), before \(L\) and
the global radius. Nonflatness therefore causes no parameter-order
circle in this Berry comparison.

This is a lemma about the finite adapted vacuum connection and the full
slow charge forms. It does not prove that all remaining full-Hamiltonian
compression and Schur terms match that covariant slow model, nor does it
prove the adapted bosonic/fermionic energy-defect estimate.
